% Letter class
\documentclass[10pt]{letter}
\usepackage[hmargin=1.5in,vmargin=1.5in]{geometry}

\usepackage[
    backend=biber,
    style=nature,
    % citestyle=authortitle,
    % natbib=true,
    % url=false,
    % doi=false,
    % eprint=false,
    % natbib=true,
    sorting=none,
]{biblatex}
\addbibresource{references.bib}

% Name of sender
\name{Thomas C. Fraser}

% Signature of sender
\signature{Thomas C. Fraser}

% Address of sender
\address
{
    \textbf{Thomas C. Fraser} \\
    tfraser@perimeterinstitute.ca\\
    +1 (226) 868-0557\\
    Perimeter Institute for Theoretical Physics,\\
    31 Caroline St N, Waterloo, Ontario N2L 2Y5, Canada
}


%-----------------------------------------------------------------------------%

\begin{document}

% Name and address of receiver
\begin{letter}
{
    Matthias Christandl \& The Hiring Committee,\\
    Centre for the Mathematics of Quantum Theory (QMATH),\\ 
    University of Copenhagen.
}

% Opening statement
\opening{Dear Professor Matthias Christandl,}

I am writing to express my interest in joining your research group as a post-doctoral student through either of the open projects ``The Mathematics of Quantum Theory'' or ``Resource Q''.
I am a PhD Candidate working under the supervision of Rob Spekkens in the Quantum Foundations group at the Perimeter Institute for Theoretical Physics in Waterloo, Ontario, Canada. I will be graduating April of 2023.

The quantum marginal problem (QMP) aims to understand the relationships between the various parts of a quantum system. Unsurprisingly, the QMP has connections to a numerous topics in quantum information theory including quantum entanglement, quantum many-body physics, quantum error correction, quantum entropic inequalities, random matrix theory, and representation theory. 

The particular route through which I discovered and became interested in tackling the QMP, while rather atypical, perfectly summarizes my academic journey. Six years ago, I became interested in the topic of causal modelling~\cite{pearl2009causality} which uses tools from computational algebraic geometry to understand the semi-algebraic constraints imposed on multivariate probability distributions under the assumption of some network of cause-effect influences between those variables. Sometimes, these algebraic constraints depend on whether or not the causal-mediaries are governed by classical physics or quantum physics, and demonstrating such a difference is an important prerequisite for establishing a potential quantum advantage in network scenarios. My first paper was the first to derive analytic constraints for a particular network involving three variables that witnesses this quantum-classical gap~\cite{fraser2018causal}. Currently, the leading strategy for deriving these constraints is to first convert the problem into a family of classical or quantum marginal problems and then afterwards solve the resulting marginal problems~\cite{wolfe2019inflation}. 

My major contribution to the field of classical causal modelling, which eventually became my master's thesis~\cite{fraser2020combinatorial}, was to switch from using tools from algebraic geometry to develop new tools which emphasized the combinatorics of the problem. This change of approach led me to develop new algorithms that were more computationally tractable than existing methods and furthermore capable of resolving some open conjectures in the field. During my PhD, I turned my attention toward quantum causal modelling, and the quantum marginal problems that it generates.

Over the last two decades, tremendous progress has been made toward solving the QMP for the special case of disjoint subsystems~\cite{klyachko2006quantum}. These advances were made possible by deep connections between i) the representation theories of both classical Lie groups and symmetric groups and ii) the convex geometry that emerges from the semigroup structure on the highest weight vectors which label irreps in the appropriate asymptotic limit. The particular manifestation of this connection which relates the spectra of quantum states to the isotypic components of Schur-Weyl decompositions was given an operational significance through Keyl \& Werner's work on spectral estimation~\cite{keyl2005estimating}. This correspondence was strengthened even further to make partial progress toward the overlapping QMP~\cite{christandl2018recoupling}.

Unfortunately, these breakthroughs were not readily applicable to the types of QMPs encountered in quantum causal modelling wherein i) the overlapping marginals are quite dense, and ii) the eigenvectors of the marginals play a critical role. In my latest paper~\cite{fraser2022sufficient}, I managed to leverage existing techniques to overcome these obstacles to derive a complete hierarchy of inequality constraints which exactly characterize the set of compatible quantum marginals. The primary tool which led to this breakthrough was to understand these asymptotic correspondences from the perspective of large deviation theory~\cite{dupuis2011weak}, from which it became clear how to freely compose these estimation schemes (in the language of Keyl~\cite{keyl2006quantum}) independently of whether or not they commute.

During my PhD, I also worked with D. Schmid, et. al.~\cite{schmid2020understanding} on the subject of entanglement theory from the perspective of abstract resource theories~\cite{coecke2016mathematical}. In that paper we emphasized the role that causal networks play when deciding which operations should be considered free, which in turn impacts how one quantifies entanglement.

% conclusion
In conclusion, I believe the tools I've learned, namely representation theory, large deviation theory, resource theory, algebraic geometry and causal modelling, in addition to the progress I've made using these tools make me an excellent addition to your research group. I'm self-driven, but eager to collaborate with others. Attached is my CV which includes a list of my papers as well as contact information for two referees. Thank you for your time and consideration.

% Closing statement
\closing{Sincerely,}

\end{letter}

\printbibliography{}
\end{document}
